%% Question bank chapter file - generated by bank_tools.py build (do not edit while the bank is still being built:
%% it is regenerated from the per-test data in Question Bank\data). Every question is one bankq block:
%%   \begin{bankq}[id=..., source={...}, paper=.., qno=.., type=..., subtopic={...}, answer={...}, answer_src=key|solved,
%%                 working={...}, difficulty=1-10, diff_parts={K S I C F}, diff_note={...},
%%                 group=<passage id or empty>, also_in={duplicate copies in other tests}]
%% difficulty = min(10, K+S+I+C+F-1): K concepts 1-3, S steps 1-4, I insight 0-2, C computation 0-2, F format 0-1
%%   ...question LaTeX, options, figures (\bfig / \bfigi -> img/)...
%%   \end{bankq}
%% Passages shared by several questions are bankpassage blocks; their questions carry group=<passage id>.
\documentclass[10pt]{article}
\usepackage[a4paper,margin=1.8cm]{geometry}
\usepackage[T1]{fontenc}
\usepackage{amsmath,amssymb,mathtools}
\usepackage{graphicx}
\usepackage[export]{adjustbox}
\usepackage{xcolor,needspace,keyval}
\usepackage{enumitem}
\usepackage{booktabs,array,tabularx}
\usepackage[hidelinks]{hyperref}
\setlength{\parindent}{0pt}
\setlength{\parskip}{3pt}
\makeatletter
\define@key{bq}{id}{\def\bq@id{#1}}\define@key{bq}{source}{\def\bq@src{#1}}\define@key{bq}{paper}{\def\bq@paper{#1}}
\define@key{bq}{qno}{\def\bq@qno{#1}}\define@key{bq}{type}{\def\bq@type{#1}}\define@key{bq}{subtopic}{\def\bq@sub{#1}}
\define@key{bq}{answer}{\def\bq@ans{#1}}\define@key{bq}{answer_src}{\def\bq@asrc{#1}}\define@key{bq}{working}{\def\bq@work{#1}}
\define@key{bq}{group}{\def\bq@group{#1}}\define@key{bq}{also_in}{\def\bq@also{#1}}
\define@key{bq}{difficulty}{\def\bq@diff{#1}}\define@key{bq}{diff_parts}{\def\bq@dparts{#1}}\define@key{bq}{diff_note}{\def\bq@dnote{#1}}
\newenvironment{bankq}[1][]{\def\bq@work{}\def\bq@also{}\def\bq@group{}\def\bq@diff{?}\def\bq@dparts{}\def\bq@dnote{}\setkeys{bq}{#1}%
  \par\needspace{4\baselineskip}\medskip\noindent{\small\textbf{\bq@id}\hfill\color{gray}[\bq@src, Paper \bq@paper, Q\bq@qno]
  \bq@type\ $\cdot$ \bq@sub\ $\cdot$ \textbf{D\bq@diff}}\par\nopagebreak\relax}
  {\par{\small\color{teal}Answer: \textbf{\bq@ans}\ifx\bq@work\@empty\else\ (solved: \bq@work)\fi
  \ifx\bq@also\@empty\else\par\color{gray}Also in: \bq@also\fi}\par}
\define@key{bp}{id}{\def\bp@id{#1}}\define@key{bp}{source}{\def\bp@src{#1}}\define@key{bp}{questions}{\def\bp@qs{#1}}
\newenvironment{bankpassage}[1][]{\setkeys{bp}{#1}\par\needspace{4\baselineskip}\bigskip\noindent
  {\small\textbf{Passage \bp@id}\hfill\color{gray}[\bp@src; questions \bp@qs]}\par\nopagebreak\relax\itshape}{\par}
\makeatother
\newcommand{\bfig}[2]{\par\smallskip{\centering\includegraphics[width=#1pt,max width=0.62\linewidth,max height=0.32\textheight]{img/#2}\par}\smallskip}
\newcommand{\bfigi}[2]{\includegraphics[width=#1pt,max width=\linewidth,max height=0.32\textheight,valign=c]{img/#2}}
\newcommand{\bmissing}[1]{\textit{[figure missing in the source conversion: \detokenize{#1}]}}
\newenvironment{qoptions}{\begin{itemize}[leftmargin=1.8em,itemsep=1pt,topsep=2pt,label={}]}{\end{itemize}}
\newcommand{\opt}[2]{\item[(#1)] #2}
\providecommand{\fwstart}{}\providecommand{\fwend}{}\providecommand{\qfullsplit}{}
\newcolumntype{L}[1]{>{\raggedright\arraybackslash}p{#1}}
\newcolumntype{C}[1]{>{\centering\arraybackslash}p{#1}}
\newcolumntype{R}[1]{>{\raggedleft\arraybackslash}p{#1}}
\providecommand{\sectionhdr}[1]{}\providecommand{\subsectionhdr}[1]{}


\begin{document}
{\Large\bfseries Units, Errors \& Experiments --- Question Bank}\par\smallskip
{\small Physics chapter P25. Questions from all converted test papers, grouped by source series. Answers: printed key where it matches the paper, otherwise solved (answer\_src=solved, with working). Exact repeats across papers appear once, with also\_in listing the other copies.}\par\medskip\hrule

\section*{TOAS papers}
\begin{bankq}[id=TOAS13-P2-PHY-S1-Q12, source={TOAS - Test 13 (24 Jul 2022)}, paper=2, qno=12, type=single, subtopic={Experiments}, answer={A}, answer_src=key, difficulty=1, diff_parts={K1 S1 I0 C0 F0}, diff_note={concepts: resistor colour code | steps: digits and multiplier from the first bands; tolerance from the last band}]
A resistance is shown in the figure. Its value (in $\Omega$) and tolerance are given respectively.
\bfig{94.2}{TOAS13_p2_phy_s1_q12_1.png}
\begin{qoptions}
\opt{A}{$14\times10^5$, 5 \%}
\opt{B}{$14\times10^5$, 10 \%}
\opt{C}{$27\times10^5$, 10\%}
\opt{D}{$27\times10^5$, 20\%}
\end{qoptions}
\end{bankq}
\begin{bankq}[id=TOAS18-P1-PHY-S1-Q03, source={TOAS - Test 18 (30 Oct 2022)}, paper=1, qno=3, type=multi, subtopic={Errors and measurement}, answer={A}, answer_src=solved, working={$\tan\theta=\omega^2(r+\ell\sin\theta)/g$ gives $g=49\times0.15/0.75=9.8$ m/s$^2$; relative error $2\Delta\omega/\omega+\Delta(\ell\sin\theta)/(r+\ell\sin\theta)+\Delta\tan\theta/\tan\theta\approx2+5.8+3.5\approx11$ percent, so C and D fail}, difficulty=6, diff_parts={K2 S2 I1 C1 F1}, diff_note={concepts: conical pendulum, propagation of errors | steps: equation for g with the rod radius; mean value; relative error of each variable; total error}]
A rod of length $r=3$ cm (exactly known) is pivoted about it's end and free to rotate in horizontal plane. A string of length $\ell=(20\pm1)$ cm is connected to opposite end of rod. A stone of mass $m$ is connected to other end of string. When rod rotate at rate $\omega=(7\pm0.07)$ rad/s then string makes angle $\theta=\left(37\pm\dfrac3\pi\right)^\circ$ with vertical. Then :
\bfig{129.4}{TOAS18_p1_phy_s1_q03_1.png}
\begin{qoptions}
\opt{A}{mean value of acceleration due to gravity is 9.8 m/s$^2$}
\opt{B}{mean value of acceleration due to gravity is 12.25 m/s$^2$}
\opt{C}{\% error in measurement of $g$ is nearly 7.67\%}
\opt{D}{\% error in measurement of $g$ is nearly 5\%}
\end{qoptions}
\end{bankq}
\begin{bankq}[id=TOAS18-P1-PHY-S3-Q05, source={TOAS - Test 18 (30 Oct 2022)}, paper=1, qno=5, type=numerical, subtopic={Errors and measurement}, answer={1.42}, answer_src=solved, working={$x=R(100-l)/l$: $\Delta x/x=\Delta R/R+\Delta l/l+\Delta l/(100-l)=1+0.25+0.167=1.42$ percent}, difficulty=2, diff_parts={K1 S1 I0 C0 F1}, diff_note={concepts: error in a meter bridge | steps: formula for the unknown resistance; sum of relative errors}]
In meter bridge setup balanced length is measured by a meter scale having least count of 0.1 cm. When known resistance of $10\ \Omega\pm0.1\ \Omega$ is using then balanced lengths are found to be 40.0 cm. Find \% error in measurement of unknown resistance $x$.
\bfig{161.1}{TOAS18_p1_phy_s3_q05_1.png}
\end{bankq}
\end{document}
